\documentclass[a4paper]{article}

%------------------------------------------------
% الحزم
%------------------------------------------------

\usepackage{enumitem}
\usepackage[margin=2.5cm]{geometry}
\usepackage{amsmath,amsfonts,amssymb}
\usepackage[charter]{mathdesign}
\usepackage{graphicx}
\usepackage{array}
\usepackage{fancyhdr}
\usepackage{polyglossia}

%------------------------------------------------
% اللغات والخطوط
%------------------------------------------------

\setdefaultlanguage[
    calendar=gregorian,
    numerals=maghrib
]{arabic}

\setotherlanguage{english}

% إذا لم تكن الخطوط مثبتة في الجهاز، تُؤخذ من المجلد ../fonts/
\IfFontExistsTF{Amiri}{%
    \newfontfamily\arabicfont[Script=Arabic, Scale=1.15]{Amiri}
}{%
    \newfontfamily\arabicfont[
        Script=Arabic, Scale=1.15,
        Path=../fonts/, Extension=.ttf,
        UprightFont=*-Regular, BoldFont=*-Bold,
        ItalicFont=*-Italic, BoldItalicFont=*-BoldItalic
    ]{Amiri}
}

\IfFontExistsTF{AlBattar}{%
    \newfontfamily\arabicfontsf[Script=Arabic, Scale=1.2]{AlBattar}
    % خط العناوين (يعمل حتى بعد استعمال \textenglish)
    \newfontfamily\khatt[Script=Arabic, Scale=1.2]{AlBattar}
}{%
    \newfontfamily\arabicfontsf[
        Script=Arabic, Scale=1.2, Path=../fonts/,
        BoldFont=ae_AlBattar.ttf
    ]{ae_AlBattar.ttf}
    % خط العناوين (يعمل حتى بعد استعمال \textenglish)
    \newfontfamily\khatt[
        Script=Arabic, Scale=1.2, Path=../fonts/,
        BoldFont=ae_AlBattar.ttf
    ]{ae_AlBattar.ttf}
}

%------------------------------------------------
% إعداد الصفحة
%------------------------------------------------

\renewcommand{\footrulewidth}{1pt}
\renewcommand{\headrulewidth}{1pt}

\rhead{\khatt المدرسة العليا للأساتذة ببوزريعة}
\lhead{\khatt ملحقة بومرداس}

\cfoot{}
\lfoot{\khatt الأستاذ: د. بن عبد الرزاق يوسف}

\parindent 0pt
\setlength{\headsep}{15pt}
\setlength{\parskip}{3pt}

\pagestyle{fancy}

\setlist[enumerate,1]{
    label=\arabic*),
    itemsep=4pt,
    topsep=3pt
}

\setlist[itemize]{
    itemsep=2pt,
    topsep=2pt
}

\arraycolsep=3pt

%================================================
% بداية الوثيقة
%================================================

\begin{document}

%------------------------------------------------
% العنوان
%------------------------------------------------



\rule{\textwidth}{1pt}

\begin{center}

{\khatt\textbf{
الفصل 4: الأخلاقيات في الذكاء الاصطناعي
}}

\end{center}

\rule{\textwidth}{1pt}

%================================================
% الوضعية المحفزة
%================================================



%================================================
% النشاط الأول
%================================================

\section*{\khatt\underline{النشاط الأول: المسؤولية والشفافية}}

\begin{enumerate}

%------------------------------------------------

\item أكمل التعريفات التالية:

\begin{center}

\begin{tabular}{|
>{\raggedleft\arraybackslash}m{4.5cm}|
>{\raggedleft\arraybackslash}m{9.5cm}|}

\hline

\textbf{المفهوم} & \textbf{التعريف} \\

\hline

المسؤولية
&
تحمل المستخدم مسؤولية استعمال الأداة والنتائج التي يعتمد
عليها.
\\
\hline

الشفافية
&
توضيح كيفية استعمال الذكاء الاصطناعي والمعلومات المتعلقة
بالمخرجات عند الحاجة.
\\
\hline

الإفصاح
&
التصريح باستعمال أداة الذكاء الاصطناعي عندما يكون ذلك
مطلوبًا أو مناسبًا.
\\
\hline

\end{tabular}

\end{center}

%------------------------------------------------

\item طالب استعمل الذكاء الاصطناعي لإنشاء ملخص لمقال علمي،
ثم قدمه للأستاذ دون الإشارة إلى أنه استعان بالأداة.

\begin{enumerate}

    \item ما المشكلة الأخلاقية في هذه الحالة؟
    
    \item هل يتحمل الطالب مسؤولية المحتوى الذي قدمه؟
    
    \item ما الإجراء الأكثر شفافية الذي يمكن للطالب القيام به؟

\end{enumerate}

%------------------------------------------------

\item أمام الحالات التالية، حدّد هل تم احترام
\textbf{الشفافية والإفصاح} أم لا:

\begin{center}

\begin{tabular}{|
>{\raggedleft\arraybackslash}m{9.5cm}|
>{\centering\arraybackslash}m{4cm}|}

\hline

\textbf{الحالة} & \textbf{التقييم} \\

\hline

الطالب يصرح بأنه استعمل أداة ذكاء اصطناعي
للمساعدة في تنظيم أفكاره.
&
\\[0.4cm]
\hline

الطالب ينسخ إجابة الذكاء الاصطناعي ويقدمها على أنها
عمله الشخصي بالكامل.
&
\\[0.4cm]
\hline

الطالب يستعمل الذكاء الاصطناعي ثم يتحقق من المعلومات
ويذكر طريقة استعماله للأداة.
&
\\[0.4cm]
\hline

الطالب يعتمد على معلومات الأداة دون التحقق منها.
&
\\[0.4cm]
\hline

\end{tabular}

\end{center}

%------------------------------------------------

\item لماذا لا تعني الاستعانة بالذكاء الاصطناعي
التخلي عن المسؤولية الشخصية؟

\end{enumerate}

%================================================
% النشاط الثاني
%================================================
\vspace{3cm}
\section*{\khatt\underline{النشاط الثاني (واجب منزلي): الانتحال القائم على الذكاء الاصطناعي}}

\textbf{يُنجز هذا النشاط في المنزل ويُسلَّم في الحصة القادمة.}

\begin{enumerate}

\item عرّف باختصار:

\begin{itemize}

    \item \textbf{الانتحال العلمي}


    \item \textbf{\textenglish{AI-Plagiarism}}
   
\end{itemize}

%------------------------------------------------

\item طالب نسخ نصًا كاملاً من أداة ذكاء اصطناعي،
ثم وضع اسمه عليه وقدمه كبحث شخصي.

\begin{enumerate}

    \item هل يعتبر هذا استعمالًا مسؤولًا؟
    
    \item ما المشكلة المتعلقة بملكية العمل؟
    
    \item ماذا كان يمكن للطالب أن يفعل بدل النسخ المباشر؟

\end{enumerate}

%------------------------------------------------

\item قارن بين الحالات التالية:

\begin{center}

\begin{tabular}{|
>{\raggedleft\arraybackslash}m{7cm}|
>{\raggedleft\arraybackslash}m{7cm}|}

\hline

\textbf{الاستعمال} & \textbf{الموقف الأخلاقي} \\

\hline

استعمال الذكاء الاصطناعي لفهم مفهوم صعب.
&
\\[0.45cm]
\hline

نسخ إجابة كاملة وتقديمها باسم الطالب.
&
\\[0.45cm]
\hline

استعمال الأداة لاقتراح أفكار ثم كتابة الطالب
للإجابة بنفسه.
&
\\[0.45cm]
\hline

استعمال الذكاء الاصطناعي لإنشاء مسودة ثم مراجعتها
والتحقق منها.
&
\\[0.45cm]
\hline

\end{tabular}

\end{center}

%------------------------------------------------

\item اقترح ثلاث قواعد للاستعمال الأخلاقي للذكاء الاصطناعي
في إنجاز البحوث الجامعية.

\end{enumerate}


\end{document}